\section{实验结果与分析}
\label{sec:results}

\subsection{评价设置与比较范围}

主要指标为 56 个风场等权的容量因子平均绝对误差（CF MAE），数值越小越好。2023 年用于验证与模型选择，2024 年用于跨年复评。下列表格报告十个随机种子的单模型指标均值，不是模型集成的误差。对成对比较定义
\begin{equation}
 \Delta(A,B)=\frac{1}{S}\sum_{s=1}^{S}
 \left[\operatorname{MAE}(A_s)-\operatorname{MAE}(B_s)\right],\qquad S=10.
\end{equation}
负差分表示 $A$ 优于 $B$。配对置信区间沿用对应实验汇总，主要刻画固定数据上的训练随机性，不代替对未见年份、未见风场或完整模型选择过程的不确定性估计。

实验分为三个层次：规范化归因协议（下称 canonical）用于辨别天气、历史和轨迹条件的作用；共同协议迁移实验用于检验条件修正能否为 P0I 骨干带来增量；最终骨干上的条件入口与摘要实验用于确定当前实现。三者的模型输入路径与训练实现并非完全一致，因而分别解释组内比较，不将不同协议的绝对误差拼接为统一排名。表~\ref{tab:main-comparison} 为当前可直接汇总的最终骨干结构对比，不包含未经统一协议评估的外部命名方法。

\subsection{最终骨干上的总体结构对比}

表~\ref{tab:main-comparison} 比较无额外摘要、不同固定读出及 MLP 摘要结构。所有配置均保留完整历史解码记忆和显式基准轨迹条件，仅改变额外摘要入口的处理。因此，CB 表示“无额外摘要”，并非“无历史信息”。参数列只描述摘要实现相对 CBH 的增量，不是全模型参数量。

\begin{table}[tbp]
 \centering
 \caption{最终骨干上的主要结构对比。指标为十种子平均 CF MAE；粗体为本表最低值。各模型均保留完整历史记忆与基准轨迹条件。}
 \label{tab:main-comparison}
 \small
 \setlength{\tabcolsep}{4pt}
 \begin{tabularx}{\linewidth}{lXrrr}
 \toprule
 配置 & 摘要输入或读出 & 2023 & 2024 & 新增参数 \\
 \midrule
 CB & 中和额外摘要入口 & 0.083964 & 0.081797 & 0 \\
 CBH & 全部历史 token 均值 & 0.083814 & 0.081507 & 0 \\
 SRECENT & 最近四个 token 均值 & 0.083672 & 0.081418 & 0 \\
 STREND & 均值与历史 token 趋势组合 & 0.083723 & 0.081452 & 0 \\
 MFLAT & 完整记忆展平后输入 MLP & 0.083739 & 0.081417 & 102,656 \\
 MPOOL & 重复均值后输入同形状 MLP & 0.083750 & 0.081436 & 102,656 \\
 HMLP & 原始历史功率输入 MLP & 0.083893 & 0.081580 & 5,888 \\
 \textbf{SLAST} & \textbf{最后一个上下文化 token} & \textbf{0.083626} & \textbf{0.081351} & \textbf{0} \\
 \bottomrule
 \end{tabularx}
\end{table}

SLAST 在本组配置中取得最低平均误差，2023/2024 年分别为 $0.083626/0.081351$。相对无额外摘要的 CB，误差降低约 $0.40\%/0.55\%$；相对均值摘要 CBH，降低约 $0.22\%/0.19\%$。这些改善幅度较小，但读出不增加参数，因而支持采用 SLAST 作为当前实现。该结论限定于本组比较，不扩展为所有架构上的最优性或普适领先性。

大型摘要网络并未形成明显成本收益优势。MFLAT 相对同容量 MPOOL 的配对差分仅为 $-0.000011/-0.000019$，两年置信区间均包含零，见表~\ref{tab:summary-paired}。这一对照尚未确认完整 token 展平输入优于均值输入。由于各配置均重新训练，摘要改变也会影响共同历史编码器的优化，不能把结果解释为固定同一历史记忆后的纯读出比较。

\subsection{条件入口与核心模块消融}

表~\ref{tab:ablation} 按协议组织关键差分。canonical 中，History 为冻结历史预测；W 增加天气修正；L 进一步让天气修正读取基准水平；A 使用严格可分离的历史、天气输出修正；J 使用联合条件修正；H 保留原始历史而中和显式轨迹；T 保留显式轨迹而移除修正侧原始历史。LP 在 L 的基础上增加基准相邻差分。JM 与 G 是用于检验乘法调制的参数匹配比较，其中 G 使用普通联合拼接。

\begin{table}[tbp]
 \centering
 \caption{关键模块消融与迁移差分。所有差分均为“左侧模型减右侧模型”，单位为 CF MAE；负值表示改善。胜出数依次为 2023/2024。“—”表示当前汇总未给出该项，不代表零胜出。不同分组只作组内解释。}
 \label{tab:ablation}
 \small
 \setlength{\tabcolsep}{4pt}
 \begin{tabularx}{\linewidth}{lXrrc}
 \toprule
 比较 & 被检验的增量 & 2023 差分 & 2024 差分 & 胜出数 \\
 \midrule
 \multicolumn{5}{l}{\textbf{canonical：信息组织与交互归因}} \\
 W$-$History & 天气修正路径 & $-0.010897$ & $-0.010383$ & 10/10；10/10 \\
 L$-$W & 显式基准水平 & $-0.007802$ & $-0.007107$ & 10/10；10/10 \\
 J$-$A & 联合修正相对严格可分离修正 & $-0.003004$ & $-0.002122$ & 10/10；10/10 \\
 J$-$H & 原始历史之外的显式轨迹 & $-0.000953$ & $-0.000838$ & 10/10；10/10 \\
 J$-$T & 显式轨迹之外的原始历史 & $-0.000918$ & $-0.001857$ & 10/10；10/10 \\
 LP$-$L & 相邻基准差分 & $-0.000035$ & $-0.000020$ & 8/10；7/10 \\
 JM$-$G & 参数匹配的乘法调制 & $+0.000014$ & $+0.000127$ & — \\
 \midrule
 \multicolumn{5}{l}{\textbf{最终骨干：条件入口与摘要}} \\
 CB$-$CNEUTRAL & 样本对应的轨迹条件 & $-0.004680$ & $-0.004442$ & 10/10；10/10 \\
 SLAST$-$CB & 末位上下文摘要入口 & $-0.000338$ & $-0.000447$ & 9/10；8/10 \\
 SLAST$-$CBH & 末位 token 相对均值读出 & $-0.000188$ & $-0.000157$ & 7/10；8/10 \\
 \midrule
 \multicolumn{5}{l}{\textbf{共同协议迁移：整体修正组织}} \\
 PJ$-$P0 & 条件修正相对原始骨干 & $-0.000367$ & $-0.000328$ & — \\
 \bottomrule
 \end{tabularx}
\end{table}

\paragraph{显式轨迹条件的作用。}
canonical 的 L$-$W 表明，天气修正显式读取基准水平具有较大预测收益。最终骨干中，CB$-$CNEUTRAL 在两年均为十种子全胜。该比较保留公共解码历史和输出基准，主要改变条件入口中的样本特异轨迹信息，因此更直接支持显式参考条件的价值。但 CB 同时包含 $b$ 和 $\Delta b$，且首步差分含当前观测功率，不能将其称为严格的“仅基准水平”对照。中和输入保留了模块形状，也不意味着有效函数自由度完全不变。

\paragraph{历史与轨迹的互补性。}
J$-$H 和 J$-$T 在两个年份均为十种子全胜，支持 canonical 修正网络中的原始历史与显式预测轨迹具有互补价值。J$-$A 则表明该联合修正优于当前严格可分离对照。然而，这些结果不能替代最终骨干中仅删除解码历史入口的实验；A 与 L 也不是嵌套对照，不能据其差异认定独立历史分支本身有害。

\paragraph{整体结构的迁移。}
共同协议下 PJ 相对 P0 的差分为 $-0.000367/-0.000328$，说明条件修正组织能够给已有骨干带来小幅增量。PJ 同时改变参考使用和条件路径，因此该差分不单独识别某一个模块。PA 使用的可分离适配器也有改善，PJ 相对 PA 未形成可靠优势；由于 PA 后续解码器仍联合处理输入，它也不等于严格输出加性模型 A。

\subsection{摘要收益及统计不确定性}

\begin{table}[tbp]
 \centering
 \caption{摘要配对比较及其原报告的 $95\%$ 置信区间。为便于阅读，差分和区间均乘以 $10^3$。数值可能与表~\ref{tab:main-comparison} 的六位小数均值相减有末位差异。}
 \label{tab:summary-paired}
 \small
 \setlength{\tabcolsep}{5pt}
 \begin{tabular}{lcc}
 \toprule
 比较 & 2023：差分 [区间] & 2024：差分 [区间] \\
 \midrule
 CBH$-$CB & $-0.150\;[-0.307,+0.028]$ & $-0.290\;[-0.517,-0.078]$ \\
 SLAST$-$CB & $-0.338\;[-0.508,-0.165]$ & $-0.447\;[-0.678,-0.234]$ \\
 SLAST$-$CBH & $-0.188\;[-0.322,-0.053]$ & $-0.157\;[-0.299,-0.003]$ \\
 MFLAT$-$CBH & $-0.076\;[-0.173,+0.026]$ & $-0.090\;[-0.277,+0.100]$ \\
 MFLAT$-$MPOOL & $-0.011\;[-0.068,+0.057]$ & $-0.019\;[-0.120,+0.094]$ \\
 \bottomrule
 \end{tabular}
\end{table}

SLAST 相对 CB 的两年区间均为负，胜出数分别为 9/10 和 8/10，支持保留额外摘要条件入口。相对均值摘要，其改进更小：两年的原配对区间均为负，但原多重比较校正后的 Holm $p$ 值为 $0.049805/0.131836$，不能称为两年均通过全部显著性标准。相对 CB 的 Holm $p$ 值为 $0.008789/0.007812$，证据强于仅比较两种摘要读出。本文依据验证集的小幅收益与零新增参数采用 SLAST，不将特定读出方式单列为方法创新。

均值摘要 CBH 相对 CB 的两年均值均改善，但 2023 年区间包含零，2024 年只有 6/10 个种子胜出。这支持“收益较小且种子间不稳定”的描述，既不能视为均值摘要必不可少，也不能将其概括为完全无效。

\subsection{参考依赖、校准解释与复杂度取舍}

已完成的分路径诊断针对\emph{均值摘要 CBH}，不将其改名为 SLAST 测量。仅替换条件路径的基准、保留输出相加路径时，供体置换使 2023/2024 MAE 约退化 $0.01369/0.01320$，而中和条件对照不受影响。这说明 CBH 依赖样本匹配的参考条件。组合路径 Jacobian 的对角均值约为 $0.9085/0.9060$，不支持在所分析样本上普遍完全抵消基准的极端解释。

上述内部替换可能制造输入组合不一致，退化幅度不能直接等同于参考的因果贡献；Jacobian 保持正向依赖也不能排除一般非线性校准。逐提前量分析表明，轨迹条件相对中和条件的改善主要在中远期；canonical 中原始历史条件的收益也随提前量增加。因此，现有结果不支持将两者简单划分为“原始历史只负责近端、预测轨迹只负责远端”。

LP$-$L 的均值差分为 $-0.000035/-0.000020$，说明相邻基准差分最多提供很小增量；2024 年区间仍包含零。该特征成本低，可以保留，但不足以支持动态响应机制。参数匹配的 JM$-$G 为正，也未支持显式乘法调制优于普通联合拼接。当前因而保留简洁非线性适配器，不增加递推隐状态或特定物理响应算子。

完整解码历史、天气直接路径、冻结训练及各上下文字段，并不因整体模型改善就自动被逐项证明必要。尤其是完整历史记忆与额外摘要应分开归因：摘要消融已经完成，最终骨干中仅中和完整解码记忆的独立再训练对照仍不在上述结果中。本文据此将主要结论限定为显式参考条件的稳定效用与历史摘要的较小增益；不从风场 MAE 改善直接推导区域聚合误差同幅改善，也不作未完成外部公平对比的最优性声明。
